%=====================================================================
%  Chapter 1 — General Presentation and State of the Art
%=====================================================================
\chapter{General Presentation and State of the Art}
\label{chap:presentation}

\section{Introduction}

Humanoid robotics is one of the most demanding fields of modern
engineering: it combines mechanical design, real-time control, embedded
electronics, and, increasingly, artificial intelligence into a single
integrated system. Reproducing a behaviour as ordinary as human walking
requires a machine to balance on a small support area while continuously
coordinating many actuators, all under strict timing constraints. This
makes a humanoid platform an ideal testbed for studying the interaction
between low-level real-time control and high-level decision making.

This chapter introduces the project by outlining its objectives and its
significance in the field of bipedal robotics. It then discusses the
problematic aspects associated with bipedal walking, presents the state
of the art through a few representative humanoid robots, and finally
positions the \textbf{Abbes} robot --- the platform developed in this
work --- with respect to those references. The objectives pursued are:

\begin{enumerate}[itemsep=2pt,topsep=4pt]
  \item Design and assemble an 18-DOF humanoid structure from 3D-printed
        parts driven by hobby servo-motors.
  \item Implement the on-board real-time firmware that actuates the
        joints and streams inertial measurements.
  \item Build a modular ROS\,2 software stack covering the kinematic
        model, balance control, and bipedal-gait generation.
  \item Validate the robot in a physics-based simulation (digital twin)
        before deployment on hardware.
  \item Add a natural-language command layer, based on a locally hosted
        Large Language Model, that translates free-form user requests
        into a fixed and safety-checked set of robot actions.
\end{enumerate}

\section{Problematic}
\label{sec:problematic}

Bipedal walking poses a unique set of challenges that stem from the need
to maintain dynamic stability on two legs. For the platform developed
here, these challenges are tightly coupled with the software-engineering
goals of the project. The core problematic aspects are:

\begin{itemize}[itemsep=2pt,topsep=4pt]
  \item \textbf{Mechanical design.} Defining a joint distribution that
        is rich enough to produce a stable, human-like gait while
        remaining buildable with low-cost, 3D-printed parts and hobby
        servo-motors.
  \item \textbf{Control system.} Despite the need for inverse kinematics
        in controlling a walking humanoid, ensuring robustness and
        adaptability in real-time control algorithms running on a
        resource-constrained microcontroller remains a significant
        challenge.
  \item \textbf{Stability criteria.} Keeping the robot balanced during
        the various phases of walking, which requires a recognised
        stability criterion and an inertial feedback loop.
  \item \textbf{ROS\,2 integration.} Leveraging the advanced integration
        capabilities of ROS\,2 to bridge a physics-based simulation and
        the physical robot through a common middleware.
  \item \textbf{High-level interaction.} Allowing a user to command the
        robot in natural language while guaranteeing that only safe,
        well-defined motions are ever executed.
\end{itemize}

\section{State of the Art}
\label{sec:sota}

Existing humanoid robots around the world showcase the remarkable
advancements in robotics technology, demonstrating capabilities ranging
from human-like motion and interaction to replacing humans in various
industrial areas. The following sections review a few representative
platforms.

\subsection{WABIAN by Waseda}
Since 1967, Waseda University has been at the forefront of biped-robot
research. WABIAN-2 features 41 mechanical degrees of freedom, including
two 6-DoF legs, a 2-DoF waist, a 2-DoF trunk, two 7-DoF arms, two 3-DoF
hands, a 3-DoF neck, and two 1-DoF feet with passive toe joints.
Standing at 1480~mm tall and weighing 64~kg, WABIAN-2 achieves a
human-like gait characterised by stretched knees, heel-contact, and
toe-off motion. This is made possible by its foot mechanism with passive
toe joints and a 2-DoF (Roll, Yaw) waist that emulates human pelvic
movement.

\begin{figure}[ht]
  \centering
  \includegraphics[height=6cm]{robot_wabian.png}
  \caption{WABIAN by Waseda.}
  \label{fig:wabian}
\end{figure}

\subsection{ASIMO by Honda}
ASIMO is a research robot and is not available for commercial sale; it
has, however, been loaned for unique events and used by major companies
as a reception host. As of February 2009, there were over 100 ASIMO
robots in various versions, ultimately aiming to assist individuals who
are disabled, elderly, or ill, and to perform tasks that are dangerous
for humans. The fifth version of ASIMO, released in September 2011,
features enhanced capabilities including the use of all its fingers,
equipped with tactile sensors in the palm and force detectors in each
finger; combined with object-recognition technology, it can pick up a
bottle and unscrew its cap, or hold a paper cup without crushing it. One
of its main innovations is the ability to change direction smoothly
while walking, without interrupting its stride.

\begin{table}[ht]
  \centering
  \caption{Comparison of ASIMO Robot Models.}
  \label{tab:asimo}
  \footnotesize
  \setlength{\tabcolsep}{4pt}
  \renewcommand{\arraystretch}{1.25}
  \begin{tabularx}{\textwidth}{@{}llllXXcc@{}}
    \toprule
    \textbf{Project} & \textbf{Date} & \textbf{Weight} & \textbf{Length}
    & \textbf{Anatomy} & \textbf{Functions} & \textbf{DoF} &
    \textbf{DoF Head} \\
    \midrule
    P-1   & 1993--1997 & 175\,kg & 1.915\,m & Head, Arm, Hand, Leg, Feet, Trunk & Walking & -- & -- \\
    P-2   & --   & 210\,kg & 1.82\,m & -- & Walking, Climbing stairs & 30 & -- \\
    P-3   & --   & 52\,kg  & 1.2\,m  & -- & Speed walking (1.6\,km/h) & 26 & -- \\
    Asimo I   & 2000 & 54\,kg & 1.3\,m & -- & Speed walking (3\,km/h), Object detection, Face recognition & 34 & 3 \\
    Asimo II  & 2005 & 54\,kg & 1.3\,m & -- & Running (6\,km/h), Gripping objects & 34 & 3 \\
    Asimo III & 2011 & 48\,kg & 1.3\,m & -- & Human daily tasks, Collision avoidance, Dynamic path planning & 57 & 3 \\
    \bottomrule
  \end{tabularx}
\end{table}

\begin{figure}[ht]
  \centering
  \includegraphics[height=6cm]{robot_asimo.png}
  \caption{ASIMO.}
  \label{fig:asimo}
\end{figure}

\subsection{HRP Humanoid Series by AIST}
The Humanoid Robotics Project (HRP) is an initiative aimed at developing
general domestic-helper robots, funded by Japan's Ministry of Economy,
Trade and Industry (METI) and the New Energy and Industrial Technology
Development Organization (NEDO). The project is led by Kawada Industries
with support from the National Institute of Advanced Industrial Science
and Technology (AIST) and Kawasaki Heavy Industries. The HRP series,
also known as Promet, began with modifications to Honda's P3 robots and
includes features such as a tele-command system. Notably, HRP-2 can
stand up after lying flat, a capability lacking in Honda's ASIMO. The
HRP-4, a bipedal humanoid, stands 150~cm tall and weighs 39~kg.

\begin{table}[ht]
  \centering
  \caption{Comparison of HRP Robot Models.}
  \label{tab:hrp}
  \scriptsize
  \setlength{\tabcolsep}{3pt}
  \renewcommand{\arraystretch}{1.3}
  \begin{tabularx}{\textwidth}{@{}llllXXcXXc@{}}
    \toprule
    \textbf{Project} & \textbf{Date} & \textbf{Weight} & \textbf{Length}
    & \textbf{Anatomy} & \textbf{Functions} & \textbf{DoF} &
    \textbf{DoF Head} & \textbf{DoF Arm} & \textbf{DoF Leg} \\
    \midrule
    HRP-1  & 1998 & 120\,kg & 1.6\,m  & Head, Trunk, Arm, Hand, Leg, Feet & Walking / Hand operations & 30 & 2 & 14: 3 shoulders, 1 elbow, 3 wrist & 12 \\
    HRP-2  & 2003 & 58\,kg  & 1.54\,m & -- & Walking on rug, Standing, Human interaction & 30 & 2 & 12: 3 shoulders, 1 elbow, 2 wrist & 12 \\
    HRP-3  & 2007 & 65\,kg  & 1.6\,m  & -- & Speed walking 2\,km/h & 42 & 2 & 12: 3 shoulders, 1 elbow, 2 wrist & 12 \\
    HRP-4C & 2009 & 43\,kg  & 1.58\,m & -- & Anthropomorphic, Face expressions & 42 & 11: 3 neck, 8 face & 4 & 12 \\
    HRP-4  & 2010 & 39\,kg  & 1.51\,m & -- & Athleticism, Running, Operating tasks & 34 & 2 & 4 & 12 \\
    \bottomrule
  \end{tabularx}
\end{table}

\begin{figure}[ht]
  \centering
  \begin{subfigure}{0.35\textwidth}
    \centering
    \includegraphics[height=6.5cm]{robot_hrp1.png}
    \caption{HRP-1}
    \label{fig:hrp1}
  \end{subfigure}
  \hspace{1.2cm}
  \begin{subfigure}{0.35\textwidth}
    \centering
    \includegraphics[height=6.5cm]{robot_hrp2.png}
    \caption{HRP-2}
    \label{fig:hrp2}
  \end{subfigure}
  \caption{HRP humanoid robots.}
  \label{fig:hrp}
\end{figure}

\subsection{Positioning of Abbes}
\label{sec:positioning}
The platforms above achieve outstanding performance but rely on custom,
expensive hardware that is out of reach for a student project. Abbes is
positioned instead as a low-cost, fully 3D-printed, 18-DOF humanoid
built around hobby servo-motors and an ESP32 microcontroller. Its
distinguishing contribution is not raw locomotion performance but its
\emph{software}: a modern ROS\,2 control stack with a physics-based
digital twin, and a natural-language command layer driven by a locally
hosted Large Language Model --- a feature uncommon on platforms of this
scale, and the main software-engineering contribution of this work.

\section{Conclusion}
This chapter explored the fundamental principles and obstacles within
bipedal robotics. By examining both the challenges and the
state-of-the-art developments, it identified the critical focal points
for the project. The subsequent chapters delve into the specific
methodologies, design considerations, and practical implementations
required to build the robot. Table~\ref{tab:bipedal} provides a concise
overview of the key features of the robots studied above.

\begin{table}[ht]
  \centering
  \caption{Comparison of Bipedal Robots.}
  \label{tab:bipedal}
  \small
  \renewcommand{\arraystretch}{1.3}
  \begin{tabularx}{\textwidth}{@{}lccX@{}}
    \toprule
    \textbf{Robot} & \textbf{Height (mm)} & \textbf{Number of DoF} &
    \textbf{Capabilities and Features} \\
    \midrule
    WABIAN by Waseda & 1480 & 41 & Human-like gait, stretched knees, heel-contact, toe-off motion \\
    ASIMO by Honda   & Varies & Varies & Bipedal locomotion, object manipulation, human interaction \\
    HRP Series by AIST & Varies & Varies & Anthropomorphic design, various locomotion modes \\
    \bottomrule
  \end{tabularx}
\end{table}
